\ECchapter{03}{Wind Energy}{How do engineers harvest the wind?}{ECGreen}

\ECObjectives{
\begin{itemize}[leftmargin=*]
\item Explain how aerodynamic lift produces rotor torque.
\item Identify the main subsystems of a modern wind turbine.
\item Interpret a wind-turbine power curve.
\item Compare onshore, fixed-bottom offshore and floating offshore solutions.
\end{itemize}}

\begin{multicols}{2}

\begin{engineeringnews}
Modern offshore wind turbines now reach ratings above 15 MW. Floating foundations
make it possible to exploit deep-water sites where winds are stronger and more
regular, but they also introduce new challenges in mooring, maintenance and grid
connection.
\end{engineeringnews}

\begin{ecarticle}
A wind turbine converts the kinetic energy of moving air into electrical energy.
The blades are shaped like aerofoils. As air flows around a blade, a pressure
difference creates lift. The tangential component of this force produces a torque
on the rotor.

The available power in the wind is
\[
P_{\mathrm{wind}}=\frac{1}{2}\rho A v^3,
\]
where $\rho$ is air density, $A$ the swept area and $v$ wind speed. The cubic term
explains why site selection is critical: a small increase in average wind speed can
produce a large increase in annual energy output.

The rotor drives either a gearbox and a high-speed generator or a direct-drive
generator. Power electronics convert the variable-frequency output into electricity
compatible with the grid. Pitch control adjusts the blade angle, while yaw control
rotates the nacelle so that the rotor faces the wind.
\end{ecarticle}

\begin{engineeringfocus}
\textbf{Main subsystems}
\begin{enumerate}[leftmargin=*]
\item Rotor, blades and hub;
\item Main shaft and bearings;
\item Gearbox or direct-drive generator;
\item Converter and transformer;
\item Pitch and yaw actuators;
\item Tower, foundations and monitoring system.
\end{enumerate}
\end{engineeringfocus}

\begin{tcolorbox}[title=\sffamily\bfseries Did You Know?,colback=yellow!10,colframe=orange!70!black]
No turbine can capture all the power in the wind. The theoretical maximum is the
\textbf{Betz limit}: 59.3\% of the kinetic power crossing the rotor area.
\end{tcolorbox}

\begin{tcolorbox}[title=\sffamily\bfseries Key Figures,colback=blue!5,colframe=ECBlue]
\begin{tabular}{lr}
Rotor diameter & 120--250 m\\
Rated power & 3--18 MW\\
Cut-in speed & 3--4 m/s\\
Rated speed & 11--15 m/s\\
Cut-out speed & 20--25 m/s\\
Capacity factor & 30--60\%
\end{tabular}
\end{tcolorbox}

\columnbreak

\begin{tcolorbox}[title=\sffamily\bfseries Wind-turbine architecture,colback=white,colframe=ECBlue]
\centering
\begin{tikzpicture}[scale=.72,transform shape,>=Latex]
\draw[very thick,ECGreen] (0,0) -- (0,3.4);
\draw[fill=ECGray,draw=ECBlue,thick] (-.7,3.4) rectangle (1.8,4.1);
\draw[fill=ECBlue!10,draw=ECBlue,thick] (1.8,3.52) rectangle (2.8,3.98);
\draw[thick] (-.7,3.75)--(-1.45,3.75);
\filldraw[ECOrange] (-1.45,3.75) circle (.12);
\foreach \a in {0,120,240}{
  \draw[very thick,ECGreen,rotate around={\a:(-1.45,3.75)}]
  (-1.45,3.75)--(-1.45,5.0);
}
\draw[->,thick] (2.9,3.75)--(4.1,3.75) node[right]{electricity};
\node[align=center] at (.55,3.75){gearbox /\\generator};
\node[align=center] at (0,-.35){tower and\\foundation};
\node[left] at (-1.5,4.7){blades};
\node[above] at (.55,4.1){nacelle};
\end{tikzpicture}
\end{tcolorbox}

\begin{tcolorbox}[title=\sffamily\bfseries Typical power curve,colback=white,colframe=ECGreen]
\centering
\begin{tikzpicture}
\begin{axis}[
width=\linewidth,height=5.0cm,
xlabel={Wind speed (m/s)},ylabel={Power (MW)},
xmin=0,xmax=26,ymin=0,ymax=3.3,
grid=major,tick label style={font=\scriptsize},
label style={font=\scriptsize},
]
\addplot[very thick,ECBlue,smooth] table[x=wind,y=power,col sep=comma]
{data/EC03/wind_power_curve.csv};
\addplot[dashed,ECOrange] coordinates {(3,0) (3,3.2)};
\addplot[dashed,ECOrange] coordinates {(15,0) (15,3.2)};
\addplot[dashed,ECOrange] coordinates {(25,0) (25,3.2)};
\end{axis}
\end{tikzpicture}
\small Cut-in, rated and cut-out regions determine how the turbine operates safely.
\end{tcolorbox}

\begin{vocabulary}
\begin{tabularx}{\linewidth}{@{}lX@{}}
blade & pale\\
hub & moyeu\\
gearbox & multiplicateur\\
yaw & orientation de la nacelle\\
pitch & calage des pales\\
swept area & surface balayée\\
capacity factor & facteur de charge
\end{tabularx}
\end{vocabulary}

\begin{questions}
\item Why does wind power depend on the cube of wind speed?
\item What is the purpose of pitch control?
\item What is the purpose of yaw control?
\item Interpret the three regions of the power curve.
\item Compare onshore and offshore wind farms.
\end{questions}

\begin{challenge}
A coastal region plans a 500 MW wind farm. Compare an onshore solution, a
fixed-bottom offshore solution and a floating offshore solution. Consider energy
yield, foundations, maintenance, grid connection, environmental impact and public
acceptance. Present a justified recommendation.
\end{challenge}

\begin{applicationexercise}
A wind turbine has a rotor diameter of 100~m and operates in air of density
$1.225~\mathrm{kg\,m^{-3}}$ at a wind speed of 10~m/s. Using
$P=\tfrac12\rho A v^3 C_p$ with $C_p=0.44$, estimate the aerodynamic power captured by the
rotor. Compare the coefficient with the Betz limit and identify two additional conversion
losses before electricity reaches the grid.
\end{applicationexercise}

\begin{speaking}
Debate the statement: ``Large wind farms should be built mainly offshore.'' Use
technical, economic and environmental arguments.
\end{speaking}

\end{multicols}

\subsection*{Sources}
\ECsource{Global Wind Energy Council}{https://gwec.net}\quad
\ECsource{IEA Wind}{https://www.iea.org/energy-system/renewables/wind}
